\scanpage{80}
\uline{\unsure{Équation|Équations} différentielle de $\unsure{\wp|p}z$}. De façon générale, si $f$ est fonction elliptique, $f'$ est fonction elliptique, donc $f$ et $f'$ sont liés par une relation polynômiale -- \uline{toute fonction elliptique satisfait à une équation différentielle polynômiale du premier ordre}. -- Pour $\wp$, on sait que $\wp'$ est de degré $2$ sur $C\langle\wp\rangle$ ; cela peut se retrouver directement : $\wp^{\prime 2}$ est fonction elliptique paire, de degré $6$, donc polynôme en $\wp$ de degré $3$ ; on voit immédiatement \deleted{\unsure{que}} \inserted{par identification} que ce polynôme est \deleted{\unsure{de la forme}}
\[
\wp'=4\wp^3\deleted{\unsure{+}}\unsure{-}20c_1\wp-28c_2\deleted{\unsure{\wp}}
\]
Soit
\[
\boxed{\wp^{\prime 2}=4\wp^3-g_2\wp-g_3\deleted{\unsure{\wp}}}\tag*{(\unsure{7|4})}
\]
\[
(g_2=20c_1,\quad g_3=28c_2)
\]
\deleted{Le calcul (en particulier du coefficient de $\wp'$) \unsure{peut être évité}) \unsure{la relation},} en remarquant que $\frac{\wp^{\prime 2}}{(\wp-e_1)(\wp-e_2)(\wp-e_3)}$ est fonction elliptique \uline{holomorphe}, (car on vérifie qu'il n'y a pas \deleted{\unsure{de}} pôles ni à l'origine, ni aux $\unsure{\omega_i|e_i}$) donc une constante, dont on vérifie qu'elle est égale à $4$ (regarder les développements de $\wp$ et $\wp'$ à l'origine) d'où
\[
\boxed{\wp^{\prime \unsure{2|?}}=4(\wp-e_1)(\wp-e_2)(\wp-e_3)}\tag*{(5)}
\]
ce qui prouve que le polynôme en $\wp$ \unsure{cherché} est $4(\wp^3-e_1)(\wp-e_2)(\wp-e_3)$. \deleted{$\sum e_i=0$, \unsure{la différence de $\wp'$ et du second membre.}}

% archive sheet number: 79
